%% ma: L10-B13-0041
%% lop: 10
%% chuong: 5
%% bai: 13
%% dang: 10.13.4
%% loai: TN4
%% muc: TH
%% dap_an: "C"
%% nguon: "chuyen-de-cac-so-dac-trung-cua-mau-so-lieu-khong-ghep-nhom-toan-10-chuong-trinh-moi.pdf (GV. Phan Nhật Linh) (P02), Bài 02 Dạng 1 Phần I Câu 37"
%% kien_thuc: "Lớp 10 Bài 13 (trung vị)"
%% trang_thai: chinh-thuc
%% ghi_chu: "Phương án A và D gốc trùng nhau (M_e = 6), đã sửa D thành 7. Tự duyệt theo yêu cầu giáo viên 22/10/2026."
\begin{ex}
Để khảo sát kết quả môn Toán trong kì thi tuyển sinh năm vừa qua của trường A, người ta chọn một mẫu gồm $100$ học sinh. Điểm môn Toán của các học sinh được cho ở bảng tần số sau:
\begin{center}
\begin{tabular}{|c|c|c|c|c|c|c|c|c|c|c|c|}\hline
Điểm & $0$ & $1$ & $2$ & $3$ & $4$ & $5$ & $6$ & $7$ & $8$ & $9$ & $10$\\\hline
Tần số & $1$ & $1$ & $3$ & $5$ & $8$ & $13$ & $19$ & $24$ & $14$ & $10$ & $2$\\\hline
\end{tabular}
\end{center}
Trung vị của mẫu số liệu trên là
\choice{$M_e=6$}{$M_e=7{,}5$}{\True $M_e=6{,}5$}{$M_e=7$}
\loigiai{Tần số tích lũy: $1;2;5;10;18;31;50;74;\ldots$ nên giá trị thứ $50$ là $6$, thứ $51$ là $7$.

$M_e=\dfrac{6+7}2=6{,}5$. Chọn C.}
\end{ex}
